%------------------------------------------------------------------------------%
\section{Avaliando Integrais Não Elementares}

Podemos utilizar as Séries de Taylor para calcular integrais de funções
cuja integração seja muito complexa ou que a primitiva não pode ser expressa
como uma combinação finita de funções elementares.
Para ilustrar essa técnica vamos aplicá-la para uma função

\begin{example}
  Calcule a Série de Taylor para a primitiva da função
  \[
    f(x) = \sin\left(x^2\right) \;.
  \]
  \solution
  Sabemos que a função seno é igual a sua Série de Taylor para todo número real,
  isso é, podemos escrever
  \[
    \sin(\alpha) = \alpha -\frac{\alpha^3}{3!} +\frac{\alpha^5}{5!} -\frac{\alpha^7}{7!} +\cdots
  \]
  Fazendo \(\alpha = x^2 \) temos
  \[
    \sin\left(x^2\right) = x^2 -\frac{x^6}{3!} +\frac{x^{10}}{5!} -\frac{x^{14}}{7!} + \cdots
  \]
  Substituindo a função pela sua série dentro da integral e integrando termo a termo, obtemos
  \begin{align*}
    \int \sin \left(x^2\right) \,dx & = \int x^2
      - \frac{x^6}{3!}
      + \frac{x^{10}}{5!}
      - \frac{x^{14}}{7!}
      + \cdots  \,dx                                    \\[5mm]
                                    & = \int x^2 \,dx
      - \int \frac{x^6}{3!} \,dx
      + \int \frac{x^{10}}{5!} \,dx
      - \int \frac{x^{14}}{7!} \,dx
      + \cdots       \\[5mm]
                                    & = C
      + \frac{x^3}{3}
      - \frac{x^7}{7\cdot 3!}
      + \frac{x^{11}}{11\cdot 5!}
      - \frac{x^{15}}{15\cdot 7!}
      + \cdots
  \end{align*}
  onde $C$ é a constante de integração.
  Para escrevermos essa série em um somatório, colocamos $x^3$
  em evidência e escrevemos as potências em termos de $x^4$
  \begin{align*}
    \int \sin \left(x^2\right) \,dx
    & = x^3 \left(
        \frac{1}{3}
      - \frac{x^4}{7\cdot 3!}
      + \frac{x^8}{11\cdot 5!}
      - \frac{x^{12}}{15\cdot 7!}
      + \cdots
      \right) \\[3mm]
    & = x^3 \left(
        \frac{\left(-x^4\right)^0}{3}
      + \frac{(-x^4)^1}{7\cdot 3!}
      + \frac{\left(-x^4\right)^2}{11\cdot 5!}
      + \frac{\left(-x^4\right)^3}{15\cdot 7!}
      + \cdots
      \right)
  \end{align*}
  onde escolhemos $C=0$ para simplificar as expressões.
  Precisamos agora identificar as sequências que constroem os valores necessários
  \[
    \begin{array}{w{c}{2cm}w{c}{3cm}w{c}{3cm}}
      \hline
      k & 3+4k & 1+2k \\ \hline
      0 & 3    & 1    \\
      1 & 7    & 3    \\
      2 & 11   & 5    \\
      3 & 15   & 7    \\ \hline
    \end{array}
  \]
  Temos então que
  \[
    \int \sin \left(x^2\right) \,dx
    = x^3 \sum_{k=0}^\infty  \frac{\left(-x^4\right)^k}{\, (3+4k) (1+2k)! \,} \;.
  \]
\end{example}

Note que continuamos sem uma expressão fechada e finita para a integral. Porém,
mesmo a série sendo uma representação infinita, ela pode ser utilizada para aproximar
o valor da integral em qualquer ponto com a precisão que desejarmos.


%------------------------------------------------------------------------------%
